\documentclass[UTF8,a4paper,10pt]{ctexart}
\usepackage[a4paper,margin=1.75cm]{geometry}
\usepackage{amsmath,amssymb,booktabs,tabularx}
\usepackage[most]{tcolorbox}
\setlength{\parindent}{2em}\setlength{\parskip}{0.3em}
\newtcolorbox{corebox}[1]{colback=gray!8,colframe=black!70,title=\textbf{#1},boxrule=.7pt,arc=1mm}
\begin{document}
\section*{B03\quad Hessian 与二次型：二阶局部形状与正定性}
\textbf{依赖：}B01。二阶信息在驻点处先去掉一阶项，再由二次型判断局部形状。
\begin{tcolorbox}[title=母接口]
$f(x+h)=f(x)+\nabla f(x)^Th+\frac12h^TH(x)h+o(\|h\|^2)$；
驻点 $\nabla f(x)=0$ 时，局部形状由 $q(h)=h^THh$ 主导。
\end{tcolorbox}
\subsection*{1.\ 符号结构}
对称 Hessian $H$ 若正定，则 $q(h)>0$（$h\ne0$），给出严格局部极小；负定给出严格局部极大；不定则存在升降两条方向，必为鞍点。半正定/半负定只能说明二阶项未否定，不能直接推出极值。
\subsection*{2.\ 代数与几何翻译}
合同变换 $H\mapsto P^THP$ 改变坐标表示但保留惯性；特征值符号是正定性的快速证据。二维判据 $\det H>0$ 且 $f_{xx}>0$ 对应正定，$\det H<0$ 对应不定；$\det H=0$ 必须检查高阶项或特殊路径。
\subsection*{3.\ 约束与边界}
若驻点不在开域内部，不能只看 Hessian；边界、约束切向和可行路径需要转交 B04 或高数 Topic07。Hessian 是点处导数对象，二次型是局部模型中的代数函数，两者不可混称。
\subsection*{4.\ 最小例子与调用}
对 $f(x,y)=x^2+4xy+5y^2$，原点是驻点，且
\[
h^THh=2h_1^2+8h_1h_2+10h_2^2
=2(h_1+2h_2)^2+2h_2^2>0\quad(h\ne0).
\]
所以原点严格局部极小。题面出现二阶偏导表、二次型或特征值符号时，先验证驻点，再调用此符号翻译。
\begin{corebox}{检查句}
是否先验证驻点？二次型是否在所有非零方向同号？退化时是否停止二阶判别并转向高阶项/边界？
\end{corebox}
\end{document}
